% !TEX root = task-1-17-4.tex
% Editable task bank. Append new BankTask blocks at the end of the document.
\documentclass[12pt,a4paper]{article}
\usepackage{physics-textbook}
\usepackage{graphicx,needspace}
\usetikzlibrary{decorations.pathmorphing}
\newif\ifShowSolutions
\ifdefined\TutorEdition
  \ShowSolutionstrue
\else
  \ShowSolutionsfalse % Change to \ShowSolutionstrue to display solutions.
\fi
\newcounter{banktask}
\newsavebox{\BankTaskBox}
\newcommand{\BankTaskContent}[3]{%
  \noindent\textbf{\thebanktask.}\quad #2\par
  \ifShowSolutions
    \Needspace{50mm}% Keep the solution heading with its opening text or diagram.
    #3\par
    \ifdefined\TutorEdition
      \ifcsname TutorSource#1\endcsname
        \nopagebreak\smallskip{\small\csname TutorSource#1\endcsname\par}%
      \fi
    \fi
  \fi
}
\newcommand{\BankTask}[3]{%
  \refstepcounter{banktask}\par\medskip
  \sbox{\BankTaskBox}{\begin{minipage}{\linewidth}%
    \BankTaskContent{#1}{#2}{#3}\end{minipage}}%
  \ifdim\dimexpr\ht\BankTaskBox+\dp\BankTaskBox\relax>\dimexpr\textheight-20mm\relax
    \Needspace{35mm}\BankTaskContent{#1}{#2}{#3}%
  \else
    \Needspace{\dimexpr\ht\BankTaskBox+\dp\BankTaskBox\relax}%
    \noindent\usebox{\BankTaskBox}\par
  \fi
  \medskip
}
\begin{document}
\BookSubsection{1.17.4}{Математический маятник}


% Source sdamgia-632: https://phys-ege.sdamgia.ru/problem?id=632
\BankTask{sdamgia-632}{
Как изменится период малых колебаний математического маятника, если длину его нити увеличить в $4$ раза?

\noindent 1) увеличится в $4$ раза\par
\noindent 2) увеличится в $2$ раза\par
\noindent 3) уменьшится в $4$ раза\par
\noindent 4) уменьшится в $2$ раза\par
}{
\smallskip\noindent\textbf{Решение}\par
Период малых колебаний математического маятника связан с длиной его нити соотношением 

$T=2 \pi \sqrt{\frac{l}{g}} .$

 Таким образом, увеличение его нити в $4$ раза приведет к увеличению периода колебаний в $2$ раза.
Ответ: $2$.
}

% Source sdamgia-28077: https://phys-ege.sdamgia.ru/problem?id=28077
\BankTask{sdamgia-28077}{
Во сколько раз уменьшится частота малых свободных колебаний математического маятника, если длину нити увеличить в $4$ раза, а массу груза уменьшить в $5$ раз?
}{
\smallskip\noindent\textbf{Решение}\par
Период колебаний математического маятника определяется формулой: 

$T=2 \pi \sqrt{\frac{l}{g}} ,$

 из которой следует, что период колебаний не зависит от массы груза, а при увеличении длины нити в $4$ раз он увеличится в $2$ раза. Частота обратно пропорциональна периоду. Следовательно, частота уменьшится в $2$ раза.

Ответ: $2$.
}

% Source sdamgia-25236: https://phys-ege.sdamgia.ru/problem?id=25236
\BankTask{sdamgia-25236}{
Первый математический маятник совершает $30$ колебаний за $20\,\text{с}$, а второй математический маятник совершает $60$ колебаний за $80\,\text{с}$. Во сколько раз длина второго маятника больше длины первого?
}{
\smallskip\noindent\textbf{Решение}\par
Период колебаний по определению равен $T= \frac{t}{N} .$ Период колебаний математического маятника находится по формуле:

$T=2 \pi \sqrt{\frac{l}{g}} .$

Из этих двух формул получаем соотношение:

$\frac{l_2}{l_1} = ( \frac{t_2 \cdot N_1}{t_1 \cdot N_2} ) ^2 = ( \frac{80 \cdot 30}{60 \cdot 20} ) ^2 =4.$

Ответ: $4$.
}

% Source sdamgia-29073: https://phys-ege.sdamgia.ru/problem?id=29073
\BankTask{sdamgia-29073}{
В научном институте стоит в шкафу математический маятник. Представим себе, что ученые сначала измерили период $T_{1}$ колебаний этого маятника на Земле, а затем  — период $T_{2}$ колебаний этого же маятника на Луне.
Чему равно отношение $\frac{T_1}{T_2} ?$ Ускорение свободного падения на Луне равно $1{,}6\,\text{м/с}^{2}$, сопротивлением воздуха на Земле можно пренебречь.
}{
\smallskip\noindent\textbf{Решение}\par
Период колебаний математического маятника определяется формулой $T=2 \pi \sqrt{\frac{l}{g}} .$ Тогда отношение периода колебаний на Земле к периоду колебаний на Луне:

$\frac{T_1}{T_2} = \sqrt{\frac{g_2}{g_1}} = \sqrt{\frac{1{,}6}{10}} =0{,}4.$

Ответ: $0{,}4$.
}

% Source sdamgia-9079: https://phys-ege.sdamgia.ru/problem?id=9079
\BankTask{sdamgia-9079}{
Математический маятник, колеблющийся с циклической частотой $\omega$  =  $3\,\text{с}^{-1}$, в нижней точке траектории имеет ускорение, равное по модулю $a$  =  $1\,\text{м/с}^{2}$. Масса груза маятника $m$  =  $900\,\text{г}$. Чему равен запас механической энергии маятника?
}{
\smallskip\noindent\textbf{Решение}\par
Циклическая частота колебаний математического маятника равна $\omega= \sqrt{\frac{g}{l}} ,$ где $l$  — длина маятника. Выражаем $l= \frac{g}{\omega ^2} .$ В нижней точке траектории тангенциальная составляющая ускорения отсутствует, ускорение является центростремительным, оно равно $a= \frac{v ^2}{l} .$ Получаем $v ^2 =al= \frac{ag}{\omega ^2} .$
В нижней точке траектории механическая энергия маятника равна его кинетической энергии. В итоге получаем:

$E= \frac{m v ^2}{2} = \frac{mag}{2\omega ^2} = \frac{0{,}9 \cdot 1 \cdot 10}{2 \cdot 3 ^2} =0{,}5\,\text{Дж}.$

Ответ: $0{,}5$.
}

% Source sdamgia-9068: https://phys-ege.sdamgia.ru/problem?id=9068
\BankTask{sdamgia-9068}{
Период колебаний математического маятника на поверхности Земли в $1{,}2$ раза больше периода его колебаний на некоторой планете. Чему равен модуль ускорения свободного падения на этой планете? Влияние атмосферы в обоих случаях пренебрежимо мало.
}{
\smallskip\noindent\textbf{Решение}\par
Период колебаний математического маятника может быть найден по формуле $T=2 \pi \sqrt{L \over g} ,$ откуда ускорение свободного падения

$g= \frac{L ( 2 \pi ) ^2}{T ^2}$

Найдем отношение ускорения свободного падения на некоторой планете и на Земле:

$\frac{g_{\,\text{п}}}{g_{\,\text{З}}} = \frac{T_{\,\text{З}} ^2}{T_{\,\text{п}} ^2} =1{,}2 ^2 =1{,}44.$

Таким образом, ускорение свободного падения на планете равно $g_{\,\text{п}}=14{,}4\,\text{м/с} ^2 .$

Ответ: $14{,}4 \,\text{м/с} ^2 .$
}

% Source sdamgia-38062: https://phys-ege.sdamgia.ru/problem?id=38062
\BankTask{sdamgia-38062}{
Маленький груз, подвешенный на идеальной пружине, находясь в положении равновесия, увеличивает её длину на величину $l_{1}$. После сообщения небольшой по модулю скорости вдоль вертикали он начинает совершать колебания. Найдите длину $l_{2}$ нерастяжимой и невесомой нити, на которой надо подвесить этот груз, чтобы он мог совершать в вертикальной плоскости малые колебания с периодом, в два раза большим, чем период колебаний на пружине.
}{
\smallskip\noindent\textbf{Решение}\par
$1$.  Период колебаний пружинного маятника равен $T_1 = 2 \pi \sqrt{\frac{m}{k}} ,$ где $m$  — масса груза, $k$ — жёсткость пружины.
$2$.  Период малых колебаний математического маятника равен $T_2 = 2 \pi \sqrt{\frac{l}{g}}.$
$3$.  Согласно условию, должно выполняться равенство $T_2 = 2T_1.$
$4$.  В соответствии со вторым законом Ньютона, в положении равновесия у пружинного маятника сумма сил, действующих на груз вдоль вертикали, должна равняться нулю. Отсюда, с учётом закона Гука $F_{\,\text{упр}}=kl_1,$ получаем: 

$mg = kl_1,$
или $\frac{m}{k} = \frac{l_1}{g} .$

Таким образом:

$T_2=2 \pi \sqrt{\frac{l}{g}} =2T_1=2 \cdot 2 \pi \sqrt{\frac{m}{k}} =2 \pi \sqrt{\frac{4l_1}{g}} .$

$5$.  Отсюда получаем, что нужное соотношение периодов колебаний двух маятников равно при $l_2 = 4l_1.$

Ответ: $l_{2}$  =  $4l_{1}$.
}


% Source C36A4C: https://ege.fipi.ru/bank/questions.php?search=1&proj=BA1F39653304A5B041B656915DC36B38&qid=C36A4C
\BankTask{C36A4C}{
Во сколько раз уменьшится частота малых свободных колебаний математического маятника, если длину нити увеличить в $9$ раз, а массу груза уменьшить в $4$ раза?
}{
\smallskip\noindent\textbf{Решение от GPT}\par
$\nu=(2\pi)^{-1}\sqrt{g/l}$; масса не влияет. При $l_2=9l_1$ частота $\nu_2=\nu_1/3$: уменьшится в $3$ раза.
}

% Source 12C0AE: https://ege.fipi.ru/bank/questions.php?search=1&proj=BA1F39653304A5B041B656915DC36B38&qid=12C0AE
\BankTask{12C0AE}{
Небольшой груз массой $0{,}1\,\text{кг}$ подвешен на невесомой нерастяжимой нити длиной $40\,\text{см}$. В результате толчка груз стал совершать незатухающие колебания в вертикальной плоскости. В таблице приведена зависимость от времени $t$ для высоты $h$ подъёма груза относительно положения равновесия.
\begin{center}{\small\setlength{\tabcolsep}{3pt}\begin{tabular}{|c|c|c|c|c|c|c|c|c|c|}\hline $t,\,\text{с}$ & $0$ & $0{,}2$ & $0{,}4$ & $0{,}6$ & $0{,}8$ & $1{,}0$ & $1{,}2$ & $1{,}4$ & $1{,}6$ \\\hline $h,\,\text{см}$ & $0$ & $12$ & $20$ & $12$ & $0$ & $12$ & $20$ & $12$ & $0$ \\\hline\end{tabular}}\end{center}
\begin{enumerate}
\item Период колебаний груза равен $0{,}8\,\text{с}$.
\item В момент времени $0{,}4\,\text{с}$ скорость груза максимальна.
\item В промежуток времени от $0{,}2$ до $1{,}0\,\text{с}$ кинетическая энергия груза достигла минимального значения $2$ раза.
\item В момент $1{,}2\,\text{с}$ кинетическая энергия груза равна $0$.
\item Максимальная скорость груза равна $2\,\text{м}/\text{с}$.
\end{enumerate}
}{
\smallskip\noindent\textbf{Решение от GPT}\par
При $t=1{,}2\,\text{с}$ груз на максимальной высоте, поэтому $E_{\text{к}}=0$. Из сохранения энергии $mv_{\max}^2/2=mgh_{\max}$: $v_{\max}=\sqrt{2\cdot10\cdot0{,}2}=2\,\text{м}/\text{с}$. Высота повторяется через полпериода, поэтому $T=1{,}6\,\text{с}$; на интервале $0{,}2\! -\!1{,}0\,\text{с}$ минимум кинетической энергии один. Ответ: $45$.
}

% Source AEDF70: https://ege.fipi.ru/bank/questions.php?search=1&proj=BA1F39653304A5B041B656915DC36B38&qid=AEDF70
\BankTask{AEDF70}{
Железный сплошной шарик совершает малые свободные колебания на лёгкой нерастяжимой нити. Затем этот шарик заменили на сплошной алюминиевый шарик такого же диаметра. Амплитуда колебаний в обоих случаях одинакова.
Как при этом изменятся частота колебаний и максимальная кинетическая энергия шарика?
\begin{enumerate}
\item увеличится
\item уменьшится
\item не изменится
\end{enumerate}
\TaskChangeTable{Частота колебаний шарика}{Максимальная кинетическая энергия шарика}
}{
\smallskip\noindent\textbf{Решение от GPT}\par
Длина маятника одинакова, поэтому $\nu=(2\pi)^{-1}\sqrt{g/l}$ не меняется. При одинаковой амплитуде $v_{\max}$ одинакова, а масса алюминиевого шарика меньше, поэтому $E_{\text{к max}}=mv_{\max}^2/2$ уменьшается. Ответ: $32$.
}

% Source 2C5D23: https://ege.fipi.ru/bank/questions.php?search=1&proj=BA1F39653304A5B041B656915DC36B38&qid=2C5D23
\BankTask{2C5D23}{
При помощи нитяного маятника необходимо экспериментально определить ускорение свободного падения. Для этого школьник взял штатив с муфтой и лапкой, нить и стальной шарик. Какие два предмета из приведённого ниже перечня оборудования необходимо дополнительно использовать для проведения этого эксперимента?
\begin{enumerate}
\item секундомер
\item линейка
\item мензурка
\item электронные весы
\item динамометр
\end{enumerate}
}{
\smallskip\noindent\textbf{Решение от GPT}\par
При малых колебаниях \(T=2\pi\sqrt{l/g}\), откуда \(g=4\pi^2l/T^2\). Секундомером (1) измеряют время \(t\) нескольких колебаний и находят \(T=t/n\), а линейкой (2) измеряют длину маятника от точки подвеса до центра шарика. Мензурка (3), электронные весы (4) и динамометр (5) не нужны: объём, масса и вес груза в эту формулу не входят. Ответ: 12.
}

% Source 1C49C9: https://ege.fipi.ru/bank/questions.php?search=1&proj=BA1F39653304A5B041B656915DC36B38&qid=1C49C9
\BankTask{1C49C9}{
Ученику необходимо на опыте обнаружить зависимость частоты свободных колебаний нитяного маятника от длины нити. У него имеется пять маятников, характеристики которых приведены в таблице. Какие два маятника необходимо взять ученику, чтобы провести данное исследование?
\begin{center}
\begin{tabular}{|c|c|c|c|}
\hline
\shortstack{№ маятника} & \shortstack{Длина нити,\\\(\text{см}\)} & \shortstack{Объём груза,\\\(\text{см}^{3}\)} & \shortstack{Материал, из которого\\сделан груз} \\
\hline
\(1\) & \(200\) & \(30\) & алюминий \\
\hline
\(2\) & \(150\) & \(60\) & алюминий \\
\hline
\(3\) & \(100\) & \(30\) & медь \\
\hline
\(4\) & \(100\) & \(60\) & медь \\
\hline
\(5\) & \(200\) & \(60\) & алюминий \\
\hline
\end{tabular}
\end{center}
}{
\smallskip\noindent\textbf{Решение от GPT}\par
Чтобы исследовать зависимость \(\nu=\frac{1}{2\pi}\sqrt{g/l}\) от длины нити, выбирают маятники с одинаковыми грузами и разными длинами нитей. У маятников 2 и 5 грузы алюминиевые, объёмом \(60\,\text{см}^3\), а длины нитей равны соответственно \(150\) и \(200\,\text{см}\). Ответ: 25.
}

% Source 25398: https://phys-ege.sdamgia.ru/problem?id=25398
\BankTask{25398}{
Маленький незаряженный шарик, подвешенный на непроводящей нити, помещен над горизонтальной диэлектрической пластиной, равномерно заряженной положительным зарядом. Размеры пластины во много раз превышают длину нити. Опираясь на законы механики и электродинамики, объясните, как изменится частота малых свободных колебаний шарика, если ему сообщить отрицательный заряд.
}{
\smallskip\noindent\textbf{Решение}\par
При отсутствии заряда шарика частота его колебаний определяется формулой

\[v = \frac{1}{T} = \frac{1}{2 \pi} \sqrt{\frac{g}{l}} .\]


\begin{center}\begin{tikzpicture}[>=stealth,line width=.6pt]

\draw(-.36,2.97)--(.36,2.97);
\foreach \x in {-.33,-.24,...,.33}{\draw(\x,2.97)--++(.14,.14);}
\draw(0,2.97)--(.14,0) node[midway,left] {$l$};
\draw(0,2.97)--(.91,.11);
\filldraw[fill=orange!45,draw=orange!85!black] (.14,0)circle(.12) (.91,.11)circle(.12);
\draw[->](.14,0)--(2.01,.26) node[midway,below left] {$x$};
\draw[->,thick](.91,.11)--(.67,.87)node[above left]{$\vec T$};
\draw[->,thick](.88,.11)--(.88,-.65)node[below left]{$\vec F_{\text{э}}$};
\draw[->,thick](.94,.11)--(.94,-.91)node[below right]{$m\vec g$};
\fill(.91,.11)circle(.035);

\end{tikzpicture}\end{center}

Если шарик зарядить, то на него со стороны заряженной пластины будет действовать электрическая сила, направленная вертикально вниз, так как пластина имеет положительный заряд, а шарик отрицательный.
Маятник совершает малые колебания, которые являются гармоническими. Потому равнодействующая сила пропорциональна смещению $\,\text{х},$ то есть $F_x= - kx,$ и имеет период колебаний, который находится по формуле $T=2 \pi \sqrt{\frac{m}{k}} ,$ где k  — коэффициент пропорциональности.
Равнодействующая сила равна $\vec F=\vec F_{\,\text{э}} + m\vec{g} + \vec{T},$ ее проекция на ось Ох $F_x= - ( F_{\,\text{э}} + mg ) \cdot \sin \alpha .$
Для малых колебаний можно считать, что $\sin \alpha \approx \frac{x}{l} .$ Тогда 

\[F_x= - ( F_{\,\text{э}} + mg ) \cdot \frac{x}{l} = - \frac{mg + F_{\,\text{э}}}{l} x.\]

Следовательно, коэффициент пропорциональности $k= \frac{F_{\,\text{э}} + mg}{l} .$ В этом случае период колебаний равен

$T=2 \pi \sqrt{\frac{ml}{mg + F_{\,\text{э}}}} ,$ а частота 

$v = \frac{1}{2 \pi} \sqrt{\frac{mg + F_{\,\text{э}}}{ml}} = \frac{1}{2 \pi} \sqrt{\frac{g}{l} + \frac{F_{\,\text{э}}}{ml}} .$

Сравнивая частоту колебаний в обоих случаях, делаем вывод, что она увеличилась.

Ответ: увеличилась.
}

\end{document}
